\lecture{2}{Thu 11 Jan 2024 10:32}{Spectral Decomposition}

\begin{proof}[Sketch of Proof of Theroem~\ref{thm:SpectralThmFiniteDim} ]
	
	The $\impliedby$ direction is simple.

	Now we show $\implies$ .
	Assume $A$ normal. Le $\lambda$ be an eigenvalue of $A$. Let $P_\lambda$ be orthogonal projection into the $\lambda$-eigenspace $E_\lambda \subset \mathscr{H}$. 

	Note: An orthogonal projector onto $E$ is defined by $P^2 = P, P^* = P, \operatorname{Im} P = E$. We can write $\text{id}_{\mathscr{H}} = P_\lambda + Q $ where $Q$ is projection onto $E_\lambda^{\perp}$.

	Then
	\begin{align*}
		A &= (P_\lambda + Q) A (P _\lambda + Q)\\
		&= P_\lambda A P_\lambda + P_\lambda A Q + Q A P_\lambda + Q A Q \\
		&= P_\lambda A P_\lambda + Q A Q 
	\end{align*}
	Now proceed inductively on $ \operatorname{dim} \mathscr{H}$ (for finite-dimensional case).
	When $\operatorname{dim} \mathscr{H} = 0, 1$ the theorem is trivial. 

	Now consider $Q A Q | _{E_\lambda^\perp}$, we know $\operatorname{dim} E_\lambda^\perp < \operatorname{dim} \mathscr{H}$. Apply the inductive hypothesis, we know $Q A Q|_{E_\lambda^\perp}$ is unitarily diagonalizable.
	That is, there exists an orthonormal basis $\mathscr{B}_1$ of $E_\lambda^\perp$ for which $Q A Q|_{E _\lambda^\perp}$ is diagonal.

	Similarly, $P_\lambda A P_\lambda |_{E_\lambda}$ is unitarily diagonalizable (for the boundary case $E_\lambda = \mathscr{H}$, we know $A = \lambda \text{id}_{\mathscr{H}}$). Thus there exists orthonormal basis $\mathscr{B}_2$ of $E_\lambda$ for which $P_\lambda A P_\lambda |_{E_\lambda}$ is diagonalizable.

	Now take $\mathscr{B} = \mathscr{B}_1 \cup \mathscr{B}_2$. This is orthonormal and we check $A$ is diagonal on this basis. \todo{check.}
\end{proof}

\begin{example}(Hilbert Space)
	\begin{itemize}
		\item $\{ \vec{0} \} = 0 , \left\langle \vec{0} , \vec{0}  \right\rangle = 0$.
		\item $\mathbb{C}$, a one-dimensional Hilbert space, with the inner product $\left\langle z, w \right\rangle  = z^* w$.
		\item $\mathbb{C}^2$, a qubit. $\mathbb{C}^d$, a qudit. 
			Often defined as ``column vectors of complex numbers of length $d$ ''. For us, define $\mathbb{C}^d$ be the Hilbert space generated by the set of symbols $\left| 0 \right\rangle , \ldots, \left| d - 1 \right\rangle $, called \textbf{kets}. The kets form an orthonormal basis of $\mathbb{C}^d$.
	\end{itemize}
\end{example}

In other words, 
\[
	\mathbb{C}^d = \mathbb{C}\left[ \left| i \right\rangle , 0 \le i <d \right] 
\] 
such that
\[
	\left\langle i \middle| j \right\rangle = \left\langle \left| i \right\rangle , \left| j \right\rangle  \right\rangle  = \delta_{i j}
.\] 

Thus, if (assuming summation notation),
\begin{align*}
	\left| \phi \right\rangle = z_i \left| i \right\rangle , \left| \psi \right\rangle  = w_i \left| i \right\rangle 
\end{align*}
then
\[
	\left\langle \phi \middle| \psi \right\rangle 
	= w_i^* z_i
.\] 


\begin{example}(Infinite Dimensional Hilbert Space)
	\begin{itemize}
		\item 
	The space $L^2(\mathbb{R}^k) = \{f: \mathbb{R}^k \to \mathbb{C}, \int _\mathbb{R}^k |f(x)|^2 d x < \infty \} $ modulo a.e. equivalence.

	These are called \textit{wave functions}. The inner product is
	\[
		\left\langle f, g \right\rangle  = \int _{\mathbb{R}^k} f(x)^* g(x) dx
	.\] 

	Remark: We will complete this space to obtain the Dirac functions.
\item Let $S$ be a (finite) set. Define $\mathbb{C}[s]$ be the Hilbert space generated by the set $S$. Note that if $S = \{0,1, \ldots, d = 1\} $, then $\mathbb{C}[S] \cong \mathbb{C}^d$.

\item If $S$ is infinite, we can construct
	\[
		\mathbb{C}[S] = \left\{ \sum_{s \in  S} z_s \left| s \right\rangle \middle| \sum_{s \in  S} |z_s|^2 < \infty  \right\} 
	.\] 
	If $S$ is countably infinite, then $\mathbb{C}[S] \cong L^2(\mathbb{R}^k)$.
	\begin{remark}
		This isomorphism does not respect \textbf{locality}. To preserve the notion of locality, one have to put $L^2(\mathbb{R}^k)$ in a topology. Precisely \ask{verify.} need to think about a section on a fiber bundle over a manifold.
	\end{remark}

	(There is only one Hilbert space with infinite dimension.)
	\ask{Ask this in Functional Analysis class.}
	\end{itemize}
\end{example}

\begin{notation}
	\[
		\mathbb{C}[\mathbb{N}] = \ell_2
	.\] 
\end{notation}

\begin{remark}
	In the old days, quantum states are described by wave functions, in $L^2(\mathbb{R})$. We can abstract from this, and view it as superposition of different orbital configurations.

	The qudit is a continuous space, but to access information we need to \textit{measure}, which \textit{discretizes} the state.
\end{remark}

\begin{definition}[Pauli Operators]
	\label{defn:PauliOperators}
\begin{itemize}
	\item 
		$ \sigma_0 = I = 
	\begin{bmatrix}
		1 & 0 \\
		0 & 1 \\
	\end{bmatrix}
	$
\item $\sigma_1 = \sigma_X = X =
	\begin{bmatrix}
		0 &1  \\
		1 &0  \\
	\end{bmatrix}
$, the \textit{swap} gate

\item $\sigma_2 = \sigma_Y = Y = 
 	\begin{bmatrix}
		0 & -i \\
		i & 0 \\
	\end{bmatrix}
	$, \textit{composite error}

\item $\sigma_3 = \sigma_Z = Z =
 	\begin{bmatrix}
		1 & 0 \\
		0 & -1 \\
	\end{bmatrix}
	$, \textit{relative phase error}

\end{itemize}

\end{definition}
\begin{remark}
	 The $Z$ gate has no classical analogue, as it is phase-dependent.
\end{remark}

These are all Hermitian operators, and they are also unitary. 

Consider complex vector space $\mathcal{B}(\mathscr{H}) = \{\text{ linear operators } \mathscr{H} \to \mathscr{H}\} $, with dimension $(\operatorname{dim} \mathscr{H}) ^2$. 

Consider a real subspace $\mathcal{B}^{s a}\left( \mathscr{H}) \right) $ of all self-adjoint operators in $\mathcal{B}(Hscr)$. 

The Pauli operators form a basis of $\mathcal{B}^{ sa }(\mathscr{H})$.

Note: They also generate a Lie algebra.


